\documentclass[10pt,a4paper]{amsart}
\usepackage[margin=19mm]{geometry}
\usepackage{amsmath,amssymb,mathtools}
\usepackage[scr=rsfs,cal=euler]{mathalfa}
\usepackage{xfrac}
\usepackage[unicode,hidelinks,bookmarks=false]{hyperref}
\newcommand{\ie}{i.e.\ }
\newcommand{\cf}{cf.\ }
\newcommand{\resp}{respectively\ }
\newcommand{\Subsection}{\subsection}
\providecommand{\nobreakdash}{\nobreak\hbox{-}}
\setlength{\emergencystretch}{2em}
\multlinegap0pt
\usepackage{bm}
\usepackage{bbm}
\usepackage{dsfont}
\usepackage{xspace}
\usepackage{cancel}
\usepackage{mathtools}
\usepackage[shortlabels]{enumitem}
\usepackage{amsthm}
\makeatletter
\@ifundefined{thm}{\newtheorem{thm}{Theorem}}{}
\makeatother
\title[Eventually nondecreasing quasi-polynomials]{Eventually nondecreasing quasi-polynomials}
\author{Benjamin Braun, Christopher O'Neill, Antwon Park}
\date{Source: arXiv:2607.19207v1 (CC BY 4.0)}
\begin{document}
\maketitle

\noindent\emph{Source and licence.} Eventually nondecreasing quasi-polynomials. Version arXiv:2607.19207v1 (CC BY 4.0), distributed under the Creative Commons Attribution 4.0 International licence, as recorded in the arXiv licence metadata for this exact version and shown on the abstract page https://arxiv.org/abs/2607.19207v1. Full rights notice: licenses/arxiv-2607.19207-CC-BY-4.0.txt.\par\smallskip

\noindent\textit{Editorial scope.} Counted unit: the Thm whose source label is \texttt{t:endh} in the article (source0.tex lines 405--422, source label \texttt{t:endh}), delivered together with the article's own complete original proof of it (source0.tex lines 424--507, one \texttt{proof} environment of 84 lines). No mathematics is written, repaired or re-derived: statement and proof are the article's own text line for line. Required context retained uncounted: none. Every definition, display and label of those lines is retained unchanged. Omitted from this package and not redistributed: 1--404, 423--423, 508--1475. Nothing inside a retained excerpt is omitted or altered: every retained body is delivered exactly as the article's lines are written, so no omission receipt of any kind accompanies this package. Citations: the retained excerpts carry no citation command at all, so no bibliography entry of the article is transcribed and no reference is redistributed. Reference adaptation: none was needed and none was made; every reference inside the delivered unit resolves inside it, so no unresolved reference or citation is produced. Printed slips kept literal: no printed word, symbol, hypothesis or number of the counted statement or of its proof was altered. Layout and class: the article's own class is \texttt{amsart} with packages including amssymb, bm, graphicx, psfrag, xcolor, enumerate, multirow, url, hyperref, subcaption, comment, algpseudocode, mathtools, xy; the wrapper is the lane's pinned \texttt{amsart} editorial wrapper at 10pt on A4, which restates the theorem-like environments the retained text needs and adds the article's own macro definitions that the retained text needs (none), with extra wrapper packages (bm, bbm, dsfont, xspace, cancel, mathtools, enumitem). The delivered numbering is the unit's own. Assets: no retained span contains a figure environment, a graphics-inclusion command, a PDF-page inclusion, a table or a TikZ picture; no image file is copied into this package, the package declares no asset dependency and no image file is redistributed with it. Licence: version arXiv:2607.19207v1 of this article is distributed under the Creative Commons Attribution 4.0 International licence, as recorded in the arXiv licence metadata for this exact version and shown on the abstract page https://arxiv.org/abs/2607.19207v1. A full rights notice is delivered at licenses/arxiv-2607.19207-CC-BY-4.0.txt. Characters: every retained excerpt is pure ASCII, so the delivered engine, XeLaTeX with its default Latin Modern text font, reports no missing character.
\begin{thm}\label{t:endh}
We have $h \in C_\textnormal{END}$ if and only if the following hold:
\begin{enumerate}[1.]
\item 
for each $r \in [0, p-2]$, either for each $\ell \in [0, d]$ the quantity
$$
A^{(\ell)}_r(h) := \sum_{j=0}^d (-j)^\ell h_{r+1,j} - \sum_{j=0}^d (-j)^\ell h_{r,j}
$$
vanishes, or the smallest $\ell$ for which $A^{(\ell)}_r(h) \ne 0$ has $A^{(\ell)}_r(h) > 0$; and
\item 
either for each  $\ell \in [0,d]$ the quantity
$$
B^{(\ell)}(h) := \sum_{j=0}^d (-j)^\ell h_{0,j}
- \sum_{j=0}^d (-j-1)^\ell h_{p-1,j}
$$
vanishes, or the smallest $\ell$ for which $B^{(\ell)}(h) \ne 0$ has $B^{(\ell)}(h) > 0$.  
\end{enumerate}
\end{thm}

\begin{proof}
Let $e_{d,k}(x_1,\ldots,x_d)$ denote the degree~$k$ elementary symmetric polynomial on $d$~variables and let $E_{d,k} := e_{d,k}(1,2,\ldots,d)$. 
We first observe that the following equality holds for any $0\leq j\leq d$, as it will be used throughout the proof:
\begin{align*}
e_{d,k}(1-j,\ldots,d-j) 
&= \sum_{\substack{A\subseteq[d]\\ |A|=k}} \, \prod_{\alpha\in A}(-j+\alpha)
= \sum_{\substack{A\subseteq[d] \\ |A|=k}} \, \sum_{\ell=0}^k(-j)^{\ell}\sum_{\substack{B\subseteq A \\ |B|=k-\ell}} \, \prod_{b\in B}b
\\
&= \sum_{\ell=0}^k (-j)^{\ell} \!\! \sum_{\substack{B \subseteq [d] \\ |B|=k-\ell}} \, \sum_{\substack{B\subseteq A \\ |A|=k}} \, \prod_{b\in B}b
= \sum_{\ell=0}^k (-j)^{\ell} \binom{d-k+\ell}{\ell} \sum_{\substack{B\subseteq [d] \\ |B|=k-\ell}} \, \prod_{b\in B}b
\\
&= \sum_{\ell=0}^k(-j)^{\ell}\binom{d-k+\ell}{\ell}E_{d,k-\ell}.
\end{align*}
Suppose $r \in [0, p-1]$.  
Observe that
\begin{align*}
L_h(tp + r)
&= \sum_{j = 0}^d h_{r,j} \binom{t+d-j}{d}
= \frac{1}{d!} \sum_{j = 0}^d h_{r,j} \sum_{k = 0}^d e_{d,k}(1-j, \ldots, d-j) t^{d-k}
\\
&= \frac{1}{d!} \sum_{j = 0}^d h_{r,j} \sum_{k = 0}^d \bigg( \sum_{\ell = 0}^k \binom{d-k+\ell}{\ell} E_{d,k-\ell} \ (-j)^\ell \bigg) t^{d-k}
\\
&= \frac{1}{d!} \sum_{k = 0}^d \sum_{\ell = 0}^k \binom{d-k+\ell}{\ell} E_{d,k-\ell} \bigg(\sum_{j = 0}^d h_{r,j} (-j)^\ell \bigg) t^{d-k}.
\end{align*}
We first consider the case where $r \ne p-1$.
Suppose that $A^{(\ell)}_r(h)=0$ for all $\ell<m$, or equivalently that
$$
\sum_{j=0}^d (-j)^\ell h_{r+1,j}
= \sum_{j=0}^d (-j)^\ell h_{r,j}
\qquad \text{for all} \qquad
\ell < m.
$$
Therefore,
\begin{align*} 
L_h(tp + r + 1) - L_h(tp + r)
&= \frac{1}{d!} \sum_{k = 0}^d \sum_{\ell = 0}^k \binom{d-k+\ell}{\ell} E_{d,k-\ell} \bigg(\sum_{j = 0}^d (h_{r+1,j} - h_{r,j}) (-j)^\ell \bigg) t^{d-k}
\end{align*}
has degree at most $d-m$ since the innermost sum vanishes for $\ell < m$.
Further, the coefficient of $t^{d-m}$ equals
\begin{equation}\label{eq:A_leading_coeff}
	\frac{1}{d!} \binom{d}{m} \sum_{j = 0}^d (h_{r+1,j} - h_{r,j}) (-j)^m
	= \frac{1}{d!} \binom{d}{m} A^{(m)}_r(h).
\end{equation}
Therefore, the $r$-th constituent of $L_h(n+1) - L_h(n)$ is eventually nonnegative if and only if $A^{(m)}_r(h)>0$.

We next consider the case $r=p-1$, which will be proved similarly.
Suppose  $B^{(\ell)}(h)=0$ for all $\ell<m$, or equivalently, that
$$
\sum_{j=0}^d (-j)^\ell h_{0,j}
= \sum_{j=0}^d (-j-1)^\ell h_{p-1,j}
\qquad \text{for all} \qquad
\ell < m.
$$
Setting $E_{d,k}' = e_{d,k}(2, \cdots, d+1)$ and $e_{d,k} = 0$ unless $0 \le k \le d$, observe that
\begin{align*}
L_h(tp + p - 1)
&= \frac{1}{d!} \sum_{j = 0}^d h_{p-1,j} \sum_{k = 0}^d e_{d,k}(1-j, \cdots, d-j) t^{d-k}
\\
&= \frac{1}{d!} \sum_{k = 0}^d \sum_{\ell = 0}^k \binom{d-k+\ell}{\ell} E_{d,k-\ell}' \sum_{j = 0}^d h_{p-1,j} (-j-1)^\ell t^{d-k}
\\
&= \frac{1}{d!} \sum_{k = 0}^d \sum_{\ell = 0}^k \sum_{n = 0}^{k-\ell} \binom{d-k+\ell+n}{n} \binom{d-k+\ell}{\ell} E_{d,k-\ell-n} \sum_{j = 0}^d h_{p-1,j} (-j-1)^\ell t^{d-k}.
\end{align*}
Additionally,
\begin{align*}
L_h(tp + p)
&= L_h((t+1)p)
\\
&= \frac{1}{d!} \sum_{s = 0}^d \sum_{\ell = 0}^s \binom{d-s+\ell}{\ell} E_{d,s-\ell} \sum_{j = 0}^d h_{0,j} (-j)^\ell (t+1)^{d-s}
\\
&= \frac{1}{d!} \sum_{s = 0}^d \sum_{\ell = 0}^s \binom{d-s+\ell}{\ell} E_{d,s-\ell} \sum_{j = 0}^d h_{0,j} (-j)^\ell \sum_{n = 0}^{d-s} \binom{d-s}{n} t^{d-s-n}
\\
&= \frac{1}{d!} \sum_{k = 0}^d \sum_{n = 0}^k \sum_{\ell = 0}^{k-n} \binom{d-k+n+\ell}{\ell} \binom{d-k+n}{n} E_{d,k-\ell-n} \sum_{j = 0}^d h_{0,j} (-j)^\ell t^{d-k}
\end{align*}
where the substitution $k = s + n$ is used in the last step.  In the final expressions for $L_h(tp + p)$ and $L_h(tp + p - 1)$, all 4 sums may be extended to an upper limit of $d$ without introducing any nonzero summands, so the coefficient of $t^{d-k}$ in the expression for 
\begin{align*}
L_h(tp + p) - L_h(tp + p - 1)
\end{align*}
is then given by
\begin{align*}
\frac{1}{d!} \sum_{k = 0}^d \sum_{n = 0}^d \sum_{\ell = 0}^d \binom{d-k+n+\ell}{n} \binom{d-k+\ell}{\ell} E_{d,k-\ell-n} \bigg(\sum_{j = 0}^d h_{0,j}(-j)^\ell - h_{p-1,j}(-j-1)^\ell \bigg).
\end{align*}
As before, for $k < m$ the coefficient of $t^{d-k}$ is zero, and the coefficient of $t^{d-m}$ has identical sign to $B^{(m)}(h)$. 
Therefore, the remaining constituent of $L_h(n+1) - L_h(n)$ is eventually nonnegative if and only if $B^{(m)}(h) > 0$.
\end{proof}

\end{document}
